\section{结论}
\label{sec:conclusion}

本文围绕固定文本原型缺少实例外观参照这一问题，提出频谱可靠性引导的测试时语义原型校准。小波分解不直接产生最终异常图，而是调制语义 patch 证据，帮助构造逐图视觉原型；视觉原型随后以保守方式校准正常/异常文本原型，最终仍在 CLIP 语义空间中完成定位。该设计把频率线索与语义判别的职责分开，也使直接融合、纯语义自参照和边界感知可靠性能够通过受控实验分别验证。

\paragraph{局限。}
当前稿件中的实验数字仍是占位目标，不能被解释为完成结果。即使未来完整方法优于原始 AnomalyCLIP，也只有在其相对 SemanticProto 的增益集中于预期类别、且稳定性与直接融合负控共同成立时，才可以把提升归因于频谱可靠性。此外，逐图原型会受背景、异常面积和 patch 分辨率影响，逻辑异常也未被直接解决。

\paragraph{未来工作。}
将频谱可靠性引导的逐图原型构造推广到其他冻结视觉--语言编码器，并研究能够显式处理背景和逻辑异常的参照构造策略。
